\section{Full Python Code for Case Studies}
\label{app:fullcode}
%%DarkGreen
\subsection{Movie Genre Classification}
\label{app:moviecase}

\begin{lstlisting}[
  aboveskip=-0.0\baselineskip,
  belowskip=-0.0\baselineskip,
  language=Python,
  showspaces=false,
  basicstyle=\ttfamily\scriptsize,
  commentstyle=\color{gray}, 
  otherkeywords={expand, filter,execute, PANDAS_DF,group_by, join, feature_domain_range, .count, cache, unique},
  keywordstyle=\color{blue},
  caption={Full code for movie genre classification.},
  captionpos=b,
  label={lst:movidecasestudy}]

  # RDFFrames imports, graph, and prefixes
  from rdfframes.knowledge_graph import KnowledgeGraph
  from rdfframes.dataset.rdfpredicate import RDFPredicate
  from rdfframes.utils.constants import JoinType
  from rdfframes.client.http_client import HttpClientDataFormat, HttpClient
  graph = KnowledgeGraph(graph_uri='http://dbpedia.org',
                         prefixes= {'dcterms': 'http://purl.org/dc/terms/',
                                  'rdfs': 'http://www.w3.org/2000/01/rdf-schema#',
                                  'dbpprop': 'http://dbpedia.org/property/',
                                  'dbpr': 'http://dbpedia.org/resource/'}) 
  
  # RDFFrames code for creating the dataframe
  dataset = graph.feature_domain_range('dbpp:starring','movie', 'actor')
  dataset = dataset.expand('actor',[('dbpp:birthPlace', 'actor_country'),('rdfs:label', 'actor_name')])
           .expand('movie', [('rdfs:label', 'movie_name'),('dcterms:subject', 'subject'),
            ('dbpp:country', 'movie_country'),('dbpo:genre', 'genre', Optional)]).cache()
  american = dataset.filter({'actor_country':['regex(str(?actor_country),"USA")']})
  prolific = dataset.group_by(['actor']).count('movie', 'movie_count', unique=True).filter({'movie_count': ['>=100']})
  movies = american.join(prolific,'actor', OuterJoin).join(dataset, 'actor', InnerJoin)
  
  # Client and execution
  output_format = HttpClientDataFormat.PANDAS_DF
  client = HttpClient(endpoint_url=endpoint, return_format=output_format)
  df = movies.execute(client, return_format=output_format)

  # Preprocessing and preparation
  import re
  import nltk
  
  def clean(dataframe):
    for i, row in df.iterrows():
      if df.loc[i]['genre'] != None:
      value =df.at[i, 'genre']
      if re.match(regex,str(value)) is not None:
          df.at[i, 'genre'] = value.split('/')[-1]
    return dataframe

  # Remove URL from the 'genre' and convert to label keys
  df=clean(df)

  # Find the most most frequent genres
  all_genres = nltk.FreqDist(df['genre'].values)
  all_genres_df = pd.DataFrame({'genre':list(all_genres.keys()), 'Count':list(all_genres.values())})
  all_genres_df.sort_values(by=['Count'],ascending=False)
  
  # In this example, use 900 movies as a cut off for the frequent movies
  most_frequent_genres = all_genres_df[all_genres_df['Count']> 900]
  df = df[df['genre'].isin(list(most_frequent_genres['genre']))]

  # Features and factorization
  from sklearn.model_selection import train_test_split
  from sklearn.preprocessing import StandardScaler

  df= df.apply(lambda col: pd.factorize(col, sort=True)[0])
  features = ["movie_name", "actor_name", "actor_country","subject","movie_country", "subject"]
  df = df.dropna(subset=['genre'])
  x = df[features]
  y = df['genre']
  x_train, x_test, y_train, y_test = train_test_split(x, y, random_state=20)
  sc = StandardScaler()
  x_train = sc.fit_transform(x_train)
  x_test = sc.fit_transform(x_test)

  # Random Forest classifier
  from sklearn.ensemble import RandomForestClassifier

  model=RandomForestClassifier(n_estimators=100)
  model.fit(x_train,y_train)
  model.fit(x_train,y_train)
  y_pred=clf.predict(x_test)
  print("Accuracy:",metrics.accuracy_score(y_test, y_pred))
    
\end{lstlisting}
 \pagebreak
\subsection{Topic Modeling}
\label{app:topiccase}

\begin{lstlisting}[
  aboveskip=0.5\baselineskip,
  belowskip=-0.0\baselineskip,
  language=Python,
  showspaces=false,
  basicstyle=\ttfamily\scriptsize,
  commentstyle=\color{gray},
  otherkeywords={KnowledgeGraph,PANDAS_DF, execute, rdfframes,expand, filter, group_by, join, entities, count, cache, unique, select_cols,apply, lower, nltk,sklearn,TfidfVectorizer,TruncatedSVD,fit,get_feature_names,lambda},
  keywordstyle=\color{blue},
  caption={Full code for topic modeling.},
  captionpos=b,
  label={lst:topiccasestudy}]
  # RDFFrames imports, graph, prefixes, and client
  import pandas as pd
  from rdfframes.client.http_client import HttpClientDataFormat, HttpClient
  from rdfframes.knowledge_graph import KnowledgeGraph
  graph = KnowledgeGraph(
          graph_uri = 'http://dblp.l3s.de',
          prefixes = {"xsd": "http://www.w3.org/2001/XMLSchema#",
          "swrc": "http://swrc.ontoware.org/ontology#",
          "rdf": "http://www.w3.org/1999/02/22-rdf-syntax-ns#",
          "dc": "http://purl.org/dc/elements/1.1/",
          "dcterm": "http://purl.org/dc/terms/",
          "dblprc": "http://dblp.l3s.de/d2r/resource/conferences/"
      })

  output_format = HttpClientDataFormat.PANDAS_DF
  client = HttpClient(endpoint_url=endpoint, port=port,return_format=output_format)

  # RDFFrames code for creating the dataframe
  papers = graph.entities('swrc:InProceedings', paper)
  papers = papers.expand('paper',[('dc:creator', 'author'),('dcterm:issued', 'date'), ('swrc:series', 'conference'),
                        ('dc:title', 'title')]).cache()
  authors = papers.filter({'date': ['>=2005'],'conference': ['In(dblp:vldb, dblp:sigmod)']}).group_by(['author'])
                   . count('paper', 'n_papers').filter({'n_papers': '>=20', 'date': ['>=2005']})
  titles = papers.join(authors, 'author', InnerJoin).select_cols(['title'])
  df = titles.execute(client, return_format=output_format)

  # Preprocessing and cleaning
  from nltk.corpus import stopwords
  df['clean_title'] = df['title'].str.replace("[^a-zA-Z#]", " ")
  df['clean_title'] = df['clean_title'].apply(lambda x: x.lower())
  df['clean_title'] = df['clean_title'].apply(lambda x: ' '.join([w for w in str(x).split() if len(w)>3])) 
  stop_words        = stopwords.words('english')
  tokenized_doc     = df['clean_title'].apply(lambda x: x.split())
  df['clean_title'] = tokenized_doc.apply(lambda x:[item for item in x if item not in stop_words])

  # Vectorization and SVD model using the scikit-learn library
  from sklearn.feature_extraction.text import TfidfVectorizer
  from sklearn.decomposition import TruncatedSVD
  vectorizer  = TfidfVectorizer(stop_words='english', max_features= 1000, max_df = 0.5, smooth_idf=True)
  Tfidf_title = vectorizer.fit_transform(df['clean_title'])
  svd_model   = TruncatedSVD(n_components=20, algorithm='randomized',n_iter=100, random_state=122)
  svd_model.fit(Tfidf_titles)  

  # Extracting the learned topics and their keyterms
  terms = vectorizer.get_feature_names()
  for i, comp in enumerate(svd_model.components_):
      terms_comp   = zip(terms, comp)
      sorted_terms = sorted(terms_comp, key= lambda x:x[1], reverse=True)[:7]
      print_string = "Topic"+str(i)+": "
      for t in sorted_terms:
          print_string += t[0] + " "
  \end{lstlisting}

 \pagebreak
\subsection{Knowledge Graph Embedding}
\label{app:kgecase}

\vspace*{8pt}
  \begin{lstlisting}[
    aboveskip=-1.0\baselineskip,
    belowskip=-0.0\baselineskip,
    language=Python,
    showspaces=false,
    basicstyle=\ttfamily\scriptsize,
    commentstyle=\color{gray},
    otherkeywords={KnowledgeGraph, execute, rdfframes, expand, filter, join, group_by, feature_domain_range, execute, PANDAS_DF, ampligraph, train_test_split_no_unseen, ComplEx, fit,mr_score, mrr_score, evaluate_performance, expand, filter, group_by, join, entities, count, cache, unique, select_cols},
    keywordstyle=\color{blue},
    caption={Full code for knowledge graph embedding.},
    captionpos=b,
    label={lst:kgecasestudy}]
    # Get all triples where the object is a URI
  from rdfframes.knowledge_graph import KnowledgeGraph
  from rdfframes.dataset.rdfpredicate import RDFPredicate
  from rdfframes.client.http_client import HttpClientDataFormat, HttpClient    
  output_format = HttpClientDataFormat.PANDAS_DF
  client = HttpClient(endpoint_url=endpoint,
                        port=port,
                        return_format=output_format,
                        timeout=timeout,
                        default_graph_uri=default_graph_url,
                        max_rows=max_rows
                        )
  dataset = graph.feature_domain_range(s, p, o).filter({o: ['isURI']})
  df = dataset.execute(client, return_format=output_format)
    
  # Train/test split and create ComplEx model from ampligraph library
  from ampligraph.evaluation import train_test_split_no_unseen 
  triples = df.to_numpy()
  X_train, X_test = train_test_split_no_unseen(triples, test_size=10000)
  #complEx model from ampligraph library
  from ampligraph.latent_features import ComplEx
  from ampligraph.evaluation import evaluate_performance, mrr_score, hits_at_n_score
  model = ComplEx(batches_count=50,epochs=300,k=100,eta=20, optimizer='adam',optimizer_params={'lr':1e-4}, 
          loss='multiclass_nll',regularizer='LP', regularizer_params={'p':3, 'lambda':1e-5}, seed=0,verbose=True)
  model.fit(X_train)
  # Evaluate embedding model
  filter_triples = np.concatenate((X_train, X_test))
  ranks = evaluate_performance(X_test, model=model, filter_triples=filter_triples,
                                use_default_protocol=True, verbose=True)
  mr  = mr_score(ranks)
  mrr = mrr_score(ranks)
  \end{lstlisting}

  %\vspace*{-30pt}
\pagebreak

\section{Description of Queries in the Synthetic Workload}
\label{app:queries}
\vspace*{-3pt}
\hfill
\begin{table*}[th!]
  \caption{Description of the queries in the sytnthetic workload.\label{tab:def-queries}}
  \resizebox{\textwidth}{!}{%
  \begin{tabular}{|c|p{10cm}|p{4cm}|p{3cm}|} 
    \hline
    \textbf{Query} & \textbf{English Description} & \textbf{\rdfframes Operators}& \textbf{\sparql Features}\\
    \hline 
    Q1 & Get a list of films in DBpedia. For each film, return the actor, language, country, genre, story, and studio, in addition to the director, producer, and title (if available). & expand (including optional predicates) & OPTIONAL, DISTINCT\\
    \hline

    Q2 & Get a list of actors available in the DBpedia or YAGO graphs. & join (outer) between two graphs, filter & OPTIONAL, FILTER, UNION
    \\
    \hline
    Q3 & Get a list of American actors available in both the DBpedia and YAGO graphs. & join (inner) between two graphs, expand, filter &  FILTER
    \\
    \hline 
    Q4 & 
    Get the nationality, place of birth, and date of birth of each basketball player in DBpedia, in addition to the sponsor, name, and president of his team (if available). & join (left outer) between two expandable datasets, expand (including optional predicates) & OPTIONAL
    \\
    \hline 
    Q5 &  Get the players (athletes) in DBpedia and their teams, group by teams, count players, and expand the team's name. &  group\_by, count, expand & GROUP BY, COUNT, DISTINCT
    \\
    \hline
    Q6 &  For films in DBpedia that are produced by any studio in India or the United States excluding 'Eskay Movies', and that have one of the following genres (film score, soundtrack, rock music, house music, or dubstep),
    get the actor, director, producer, time and language. & expand, filter & FILTER
    \\ 
    \hline
    Q7 & For the films in DBpedia, get actors, director, country, producer, language, title, genre, story, and studio. Filter on country, studio, genre, and runtime. &  expand, filter & FILTER
    \\
    \hline
    Q8 & Get the nationality, place of birth, and date of birth of each basketball player in DBpedia, in addition to the sponsor, name, and president of his team. & join (inner) between two expandable datasets, expand & Multiple conjunctive graph patterns
    \\
    \hline 
    Q9 & Get the list of basketball players in DBpedia, their teams, and the number of players on each team. &  group\_by, count, expand & GROUP BY, COUNT, DISTINCT
    \\
    \hline
    Q10 &  For films in DBpedia that are produced by any studio in India or the United States excluding 'Eskay Movies', and that have one of the following genres (film score, soundtrack, rock music, house music, or dubstep),
    get the actor and language, in addition to the producer, director, and title (if available). & expand (including optional predicates), filter & OPTIONAL, FILTER
    \\ 
    \hline
    Q11 & Get the list of athletes in DBpedia. For each athlete, return his birthplace and the number of athletes who were born in that place. &  group\_by, count, expand & GROUP BY, COUNT, DISTINCT
    \\ 
    \hline
    Q12 & Get the pairs of films in DBpedia that belong to the same genre and are produced in the same country. For each film in each pair, return the actor, country, story, language, genre, and studio, in addition to the director, producer, and title (if available). & group\_by on multiple columns, count, expand (including optional predicates) &  GROUP BY, COUNT, OPTIONAL, DISTINCT
    \\
    \hline
    Q13 & Get the sponsor, name, president, and the number of basketball players of each basketball team in DBpedia. & join (inner) between two
    datasests (expandable, group\_by) & GROUP BY, COUNT, DISTINCT
    \\
    \hline 
    Q14 & Get the sponsor, name, president, and the number of basketball players (if available) of each basketball team in DBpedia. & join (left outer) between two datasets (expandable, group\_by), expand (including optional predicates) & GROUP BY, COUNT, OPTIONAL, DISTINCT
    \\
    \hline 
    Q15 & Get a list of the books in DBpedia that were written by American authors who wrote more than two books. For each author, return the birth place, country, and education, and for each book return the title, subject, country (if available), and publisher (if available). &  join (outer), group\_by, count, expand (including optional predicates), filter & GROUP BY, COUNT, HAVING,
    OPTIONAL, FILTER, UNION
    \\
    \hline

    Q16 & Get a list of people in the DBpedia graph who were born in the United States. Get a list of authors from the DBLP graph who have publications dated after 2015. Get a list of people in the YAGO graph who are citizens of the United States. Join the three lists retrieved from the three graphs on name. & join (full outer) between three graphs, expand (including optional predicates), filter & OPTIONAL, FILTER, DISTINCT, UNION
    \\
    \hline
  
  \end{tabular}
  }
\end{table*}

\pagebreak
%\vspace*{-30pt}
\section{Naive \sparql Query for Movie Genre Classification}
\label{app:naiveQuery}
\vspace*{-2pt}
%\hfill

\begin{lstlisting}[
  breaklines=true,
  columns=fullflexible,
  aboveskip=-0.0\baselineskip,
  belowskip=-0.0\baselineskip,
  language=SQL,
  showspaces=false,
  basicstyle=\ttfamily\scriptsize,
  commentstyle=\color{gray},
  keywordstyle=\color{blue},
  otherkeywords={OPTIONAL, FILTER},
  caption={Naive \sparql query corresponding to the \sparql query shown in Listing~\ref{lst:genresparql}.},
  captionpos=b,
  label={lst:naivemovie}
  ]
  PREFIX  dbpp: <http://dbpedia.org/property/>
  PREFIX  dcterms: <http://purl.org/dc/terms/>
  PREFIX  rdfs: <http://www.w3.org/2000/01/rdf-schema#>
  PREFIX  dbpo: <http://dbpedia.org/ontology/>
  PREFIX  dbpr: <http://dbpedia.org/resource/> 
  SELECT DISTINCT  ?actor_name ?movie_name ?actor_country ?subject ?genre
  FROM <http://dbpedia.org> WHERE
    {{{ SELECT  * WHERE
              {{ SELECT  * WHERE
                  { { SELECT  ?movie ?actor WHERE
                        { ?movie  dbpp:starring  ?actor } }
                    { SELECT  ?actor ?actor_country WHERE
                        { ?actor  dbpp:birthPlace  ?actor_country } }
                    { SELECT  ?actor ?actor_name WHERE
                      { ?actor  rdfs:label  ?actor_name } }
                    { SELECT  ?movie ?movie_name WHERE
                      { ?movie  rdfs:label  ?movie_name } }
                    { SELECT  ?movie ?subject WHERE
                        { ?movie  dcterms:subject  ?subject } }
                    { SELECT  ?movie ?movie_country WHERE
                         { ?movie  dbpp:country  ?movie_country } }
                    { SELECT  ?actor ?actor_country WHERE
                        { ?actor  dbpp:birthPlace  ?actor_country
                          FILTER regex(str(?actor_country), "USA") } }
                    { SELECT  ?movie ?genre WHERE
                        { OPTIONAL
                            { ?movie  dbpo:genre  ?genre } } } } }
              OPTIONAL
                { SELECT DISTINCT  ?actor (COUNT(DISTINCT ?movie) AS ?movie_count) WHERE
                    { { SELECT  ?movie ?actor WHERE
                          { ?movie  dbpp:starring  ?actor } }
                      { SELECT  ?actor ?actor_country WHERE
                          { ?actor  dbpp:birthPlace  ?actor_country } }
                      { SELECT  ?actor ?actor_name WHERE
                          { ?actor  rdfs:label  ?actor_name } }
                      { SELECT  ?movie ?movie_name WHERE
                          { ?movie  rdfs:label  ?movie_name } }
                      { SELECT  ?movie ?subject WHERE
                          { ?movie  dcterms:subject  ?subject } }
                      { SELECT  ?movie ?movie_country WHERE
                          { ?movie  dbpp:country  ?movie_country } }
                      { SELECT  ?movie ?genre WHERE
                          { OPTIONAL
                              { ?movie  dbpo:genre  ?genre } } } }
                  GROUP BY ?actor
                  HAVING ( COUNT(DISTINCT ?movie) >= 100 ) } } }
        UNION
          { SELECT  * WHERE
            { { SELECT DISTINCT  ?actor (COUNT(DISTINCT ?movie) AS ?movie_count) WHERE
                  { { SELECT  ?movie ?actor WHERE
                        { ?movie  dbpp:starring  ?actor } }
                    { SELECT  ?actor ?actor_country WHERE
                        { ?actor  dbpp:birthPlace  ?actor_country } }
                    { SELECT  ?actor ?actor_name WHERE
                        { ?actor  rdfs:label  ?actor_name } }
                    { SELECT  ?movie ?movie_name WHERE
                        { ?movie  rdfs:label  ?movie_name } }
                    { SELECT  ?movie ?subject WHERE
                        { ?movie  dcterms:subject  ?subject } }
                    { SELECT  ?movie ?movie_country WHERE
                        { ?movie  dbpp:country  ?movie_country } }
                    { SELECT  ?movie ?genre WHERE
                        { OPTIONAL { ?movie  dbpo:genre  ?genre } } } }
                  GROUP BY ?actor
                  HAVING ( COUNT(DISTINCT ?movie) >= 100 ) }
                OPTIONAL
                { SELECT  * WHERE
                    { { SELECT  ?movie ?actor WHERE
                          { ?movie  dbpp:starring  ?actor } }
                      { SELECT  ?actor ?actor_country WHERE
                          { ?actor  dbpp:birthPlace  ?actor_country } }
                      { SELECT  ?actor ?actor_name WHERE
                          { ?actor  rdfs:label  ?actor_name } }
                      { SELECT  ?movie ?movie_name WHERE
                          { ?movie  rdfs:label  ?movie_name } }
                      { SELECT  ?movie ?subject WHERE
                          { ?movie  dcterms:subject  ?subject } }
                      { SELECT  ?movie ?movie_country WHERE
                        { ?movie  dbpp:country  ?movie_country } }
                      { SELECT  ?actor ?actor_country WHERE
                          { ?actor  dbpp:birthPlace  ?actor_country
                            FILTER regex(str(?actor_country), "USA") } }
                      { SELECT  ?movie ?genre WHERE
                          { OPTIONAL { ?movie  dbpo:genre  ?genre } } } } } } } } }
  \end{lstlisting}

\vspace*{5pt}
%\pagebreak
\pagebreak
\section{Naive \sparql Query for Topic Modeling}
\label{app:naiveQuery2}
\hfill

\begin{lstlisting}[
breaklines=true,
columns=fullflexible,
aboveskip=-1.0\baselineskip,
belowskip=-0.0\baselineskip,
language=SQL,
showspaces=false,
basicstyle=\ttfamily\scriptsize,
commentstyle=\color{gray},
keywordstyle=\color{blue},
otherkeywords={OPTIONAL, FILTER},
caption={Naive \sparql query corresponding to the \sparql query shown in Listing~\ref{lst:topicsparql}.},
captionpos=b,
label={lst:naivetopic}
]
PREFIX  swrc: <http://swrc.ontoware.org/ontology#>
PREFIX  rdf:  <http://www.w3.org/1999/02/22-rdf-syntax-ns#>
PREFIX  owl:  <http://www.w3.org/2002/07/owl#>
PREFIX  dcterm: <http://purl.org/dc/terms/>
PREFIX  xsd:  <http://www.w3.org/2001/XMLSchema#>
PREFIX  rdfs: <http://www.w3.org/2000/01/rdf-schema#>
PREFIX  foaf: <http://xmlns.com/foaf/0.1/>
PREFIX  dc:   <http://purl.org/dc/elements/1.1/>
PREFIX  dblprc: <http://dblp.l3s.de/d2r/resource/conferences/>

SELECT  ?title
FROM <http://dblp.l3s.de>
WHERE
  { 
    {SELECT ?paper WHERE {?paper rdf:type swrc:InProceedings}}.
    {SELECT ?paper ?author WHERE {?paper dc:creator ?author}}.       
    {SELECT ?paper ?date WHERE {?paper dcterm:issued  ?date }}.   
    {SELECT ?paper ?conference WHERE {?paper swrc:series ?conference}} .
    {SELECT ?paper ?title WHERE {?paper  dc:title  ?title}}.
    {SELECT ?paper ?date WHERE {?paper dcterm:issued  ?date FILTER ( year(xsd:dateTime(?date)) >= 2005 )  }}

    { SELECT ?author WHERE
      {
       { SELECT  ?author COUNT(?paper) as  ?count_paper
          WHERE
            {
              {SELECT ?paper WHERE {?paper rdf:type swrc:InProceedings}}.
              {SELECT ?paper ?author WHERE {?paper dc:creator ?author}}.       
              {SELECT ?paper ?date WHERE {?paper dcterm:issued  ?date }}.   
              {SELECT ?paper ?conference WHERE {?paper swrc:series ?conference}} .
              {SELECT ?paper ?title WHERE {?paper  dc:title  ?title}}.
              {SELECT ?paper ?date WHERE {?paper dcterm:issued  ?date FILTER ( year(xsd:dateTime(?date)) >= 2000 )  }} .
              {SELECT ?paper ?conference WHERE {?paper swrc:series ?conference FILTER( ?conference IN (dblprc:vldb, dblprc:sigmod) )}}
            }
          GROUP BY ?author
        }
        FILTER ( ?count_paper >= 20 ) 
      }
    }
  }
\end{lstlisting}
